\paragraph{Lemma30.} Let $F$ be a system of unsatisfiable polynomials and let $z$ be a nonsingular extension variable with the extension axiom $z-Q$. Let $\ell=\ell(\supp(Q))$ so that $z^{\ell}=c$ holds for some $c\in\mathbb F_p$. Let $\Pi$ be a refutation of $F\cup\{z-Q\}$ modulo $z^{\ell}=c$ such that for at least an $\epsilon$ fraction of pairs $(t_1,t_2)$ in $H_d(\Pi)$, $\Qdeg(t_1,t_2,z)=1$, for some $d\geq0$. Then there exists a refutation $\Pi'$ of $F$ such that $|H_d(\Pi')|\leq(1-\epsilon/4p^2)|H_d(\Pi)|$.
\begin{proof}
We will apply a procedure as described by Algorithm1 in order to modify the proof to satisfy the post-conditions of the Lemma. Here we give a detailed description of the algorithm, together with its correctness. Let $H_d(\Pi,z)$ be the set of all unordered pairs $(t_1,t_2)\in H_d(\Pi)$ that $z$ contributes to. That is, $H_d(\Pi,z)$ is the set of all unordered pairs $(t_1,t_2)\in H_d(\Pi)$ such that $\Qdeg(t_1,t_2,z)=1$. There are many different ways that $z$ can contribute to $H_d(\Pi,z)$: namely, for all $i,j$ such that $i<j<\ell$, let $H_d(\Pi,z,i,j)$ be the set of all unordered pairs $(t_1,t_2)\in H_d(\Pi,z)$, such that the degree of $z$ in $t_1$ is $i$ and the degree of $z$ in $t_2$ is $j$. Note that for any two pairs $(i,j)$ and $(i',j')$ such that $i\neq i'$ or $j\neq j'$, $H_d(\Pi,z,i,j)$ and $H_d(\Pi,z,i',j')$ are disjoint. Therefore, there exists a ``good'' pair $\ell_0<\ell_1<\ell$ such that removing $H_d(\Pi,z,\ell_1,\ell_0)$ from $H_d(\Pi,z)$ will remove at least a $1/p^2$ fraction of $H_d(\Pi,z)$ and therefore a $\epsilon/p^2$ fraction of pairs in $H_d(\Pi)$, since $|H_d(\Pi,z)|\geq\epsilon|H_d(\Pi)|$.
We want to apply the Split operation $\Split_{z,\ell_0,\ell_1}$ to remove all such pairs. But in order to do this we have to satisfy the preconditions of Lemmas23 and25: we need two values $a,b$ such that $a^{\ell_1-\ell_0}\neq b^{\ell_1-\ell_0}$ and all the axioms should be free of $z$ except for $(z-a)(z-b)=0$. The first step (Line2 of1) is to apply Lemma22 with $l=\ell_1-\ell_0$. This gives us $w^*\in\vars(Q)$, $a,b\in\supp(Q)$ and a partial restriction $\delta$ to $\vars(Q)-w^*$ such that $(z-Q)|_{\delta}=z-(b-a)w^*-a$, where $a^{\ell_1-\ell_0}\neq b^{\ell_1-\ell_0}\bmod p$. Next, we apply the restriction $\delta$ to $\Pi$ (Line3).
Now we have a simpler linear extension axiom for $z$ of the form $z-(b-a)w^*-a=0$. Next we would like to make the substitution $w^*=(z-a)/(b-a)$ in $\Pi$ in order to satisfy this extension axiom, towards the goal of eliminating $z$ from the axioms so that we have the preconditions of Lemma25 and therefore are able to apply $\Split_{z,\ell_1,\ell_0}$. However, if $w^*$ appears in any of the axioms in $F$, this would create additional occurrences of $z$ and we would not make any progress. Therefore, we have to make sure that none of the axioms of $F$ contain $w^*$. But we also cannot set $w^*$ to zero or one in an attempt to get rid of it, since this would set $z$ to either $a$ or $b$ through the above extension axiom, and Split requires that $z$ take on two distinct values. We thus have to get rid of all axioms mentioning $w^*$ either by setting other variables or by replacing these axioms with stronger versions, such that the former can be derived from the latter. This is what the subroutine Cleanup does, in addition to removing the axioms in $F$ that were affected by our earlier restriction $\delta$, so that we have a clean version of $F^{\mathsf{SEL}}_{n,k}$ as defined in the previous section.
We are now ready to make the substitution $w^*=(z-a)/(b-a)$. Under this substitution, the Boolean axiom $w^{*2}-w=0$ reduces to $(z-a)(z-b)=0$, and the original extension axiom for $z$ disappears (since under this substitution it becomes $0=0$.) Thus this substitution would satisfy all of the preconditions of Lemmas23,25. However, this substitution can create a new problem: it can cause a blow up in the size of $H_d(\Pi)$ since for every pair of terms $(t_1,t_2)$ such that one of them contains $w^*$, we could have up to four new terms after the substitution. In order to deal with this potential blow up we do a simple case analysis: If $w^*$ contributes to at least an $\epsilon/4p^2$ fraction of pairs $(t_1,t_2)$ in $H_d(\Pi)$, then we set $w^*=0$ (Lines4–5). This gives us the required reduction in the size of $H_d(\Pi)$ ($z$ is also set to a constant by setting $w^*=0$, but we don't care about that since we have obtained a reduction in high Quadratic degree terms without needing to use Split). Otherwise, the blowup caused by the substitution $w^*=(z-a)/(b-a)$ adds at most $3\epsilon/4p^2$ fraction of pairs to $H_d(\Pi)$, and thus if we remove all pairs in $H_d(\Pi,z,\ell_0,\ell_1)$ (after this blowup) then overall we will will have reduced the size of $H_d(\Pi)$ to $(1-\epsilon/4p^2)|H_d(\Pi)|$. So in this latter case, we apply the substitution mentioned above (Line8) which simultaneously removes $w^*$ from all axioms, and replaces the linear axiom for $z$ by $(z-a)(z-b)=0$. Now all preconditions for Lemma8 hold so we can apply $\Split_{z,\ell_0,\ell_1}$ (Line9) to get a valid refutation. It is left to argue that this indeed removes the set $H_d(\Pi,z,\ell_1,\ell_0)$. More precisely, we argue that high Quadratic degree pairs of terms in the refutation obtained after applying Split have a one to one mapping to the set $H_d(\Pi)-H_d(\Pi,z,\ell_1,\ell_0)$. Fix a line $P\in\Pi$. Since we are working modulo $z^{\ell}=c$, we can assume that $P=P_0+zP_1+\ldots z^{\ell-1}P_{\ell-1}$. Let $R=z^{\ell_0}R_0+z^{\ell_1}R_1$ be the unique polynomial equivalent to $P\bmod(z-a)(z-b)$. $\Split_{z,\ell_0,\ell_1}(\Pi)$ is the refutation with lines $R_1,R_0$ for all $P\in\Pi$. By the proof of Lemma23 $R_0,R_1$ have the form:
\[R_1=P_{\ell_1}+\sum_{i<\ell,i\neq\ell_0}c_{1i}P_i,\qquad R_0=P_{\ell_0}+\sum_{i<\ell,i\neq\ell_1}c_{0i}P_i\]
for some constants $c_{1i},c_{0i}\in\mathbb F_p$.
For a pair of terms $(t_i,t_j)$ in $R_1$ such that $t_i\in P_i$ and $t_j\in P_j$ and $\Qdeg(t_i,t_j)\geq d$, we map it to the pair $(t_iz^i,t_jz^j)\in P$, and similarly for $R_0$. Clearly this is a one-one mapping, and since $P_{\ell_0}$ does not occur in $R_1$ and $P_{\ell_1}$ does not occur in $R_0$, it is a mapping to $H_d(\Pi)-H_d(\Pi,z,\ell_1,\ell_0)$. Therefore we have that for the refutation $\Pi'=\Split_{z,\ell_0,\ell_1}(\Pi)$ whose lines are $\{R_1,R_0\}$,
\[|H_d(\Pi')|\leq|H_d(\Pi)-H_d(\Pi,z,\ell_1,\ell_0)|\leq(1-\epsilon/4p^2)|H_d(\Pi)|.\]
\end{proof}
